\documentclass[11pt]{article}
\usepackage[T1]{fontenc}
\usepackage{amsmath}
\usepackage{amssymb}
\usepackage{mathtools}
\usepackage{booktabs}
\usepackage{longtable}
\usepackage{tabularx}
\usepackage{array}
\usepackage[margin=0.82in]{geometry}
\usepackage[colorlinks=true,linkcolor=blue,citecolor=blue,urlcolor=blue,pdfusetitle]{hyperref}

\newcommand{\R}{\mathbb{R}}
\newcommand{\norm}[1]{\left\lVert #1 \right\rVert}
\newcommand{\Evidence}[1]{\textsf{#1}}
\newcommand{\Status}[1]{\textbf{Status:} \Evidence{#1}.}
\newcommand{\dd}{\mathrm{d}}
\newcommand{\TPD}{\operatorname{TPD}}
\newcommand{\rank}{\operatorname{rank}}

\title{Formal Equilibrium Formulations for ePC-SAFT\\
\large Source, numerical, and implementation boundaries}
\author{ePC-SAFT project}
\date{13 July 2026}

\begin{document}
\maketitle
\tableofcontents

\section{Purpose, evidence discipline, and capability boundary}
\label{sec:eqform:purpose}

This document is the canonical mathematical owner for equilibrium formulations
in the M4 package context. It does not replace the EOS equation catalogue or
the algorithm inventory, and it deliberately does not import either standalone
document. It separates thermodynamic formulations that share quantities but
do not share coordinates, controllers, certificates, or globality claims.

\paragraph{Evidence vocabulary.}
\Evidence{verified} means that a statement is supported by a governing primary
source, an executable repository owner, or both, with the type of support named
nearby. \Evidence{inference} marks a mathematical consequence derived from
verified equations but not stated verbatim by the source. \Evidence{assumption}
marks a declared modelling or numerical choice needed to make a problem
well-posed. \Evidence{unknown} marks missing or ambiguous source evidence; it
must not be converted into a default, parameter, or capability claim.

\paragraph{Current public surface.}
\Status{verified implementation evidence}
The complete production-exposed equilibrium route set is
\texttt{bubble\_pressure}, \texttt{dew\_pressure}, and scoped
\texttt{single\_component\_vle}. The pure route is restricted to
nonassociating methane on \(100\)--\(180\,\mathrm{K}\), ethane on
\(100\)--\(280\,\mathrm{K}\), and propane on \(100\)--\(340\,\mathrm{K}\),
with inclusive endpoints from the retained NIST evidence. Bubble/dew
association is restricted to exact source-backed Gross--Sadowski 2002
Figures 2--9 fixtures and parameter fingerprints. Ion-containing inputs are
not public boundary inputs. None of the broader formulations below is admitted
by this document.

\begin{table}[htbp]
\centering
\small
\begin{tabularx}{\textwidth}{>{\raggedright\arraybackslash}p{0.32\textwidth}X}
\toprule
Family & Exact current implementation status \\
\midrule
Public pressure boundaries and scoped pure saturation &
production exposed only in the scope stated above \\
Neutral HELD & internal components; controller absent/deferred \\
Association-aware neutral HELD & internal components; controller absent/deferred \\
Perdomo modified-mole HELD2 & absent/deferred \\
Ascani counterion-pair equilibrium & internal components; controller absent/deferred \\
Standalone CE & public schemas/helpers; executable solver internal validation only \\
CPE & absent/deferred \\
\bottomrule
\end{tabularx}
\caption{Capability status. Descriptive family names are not activation evidence.}
\label{tab:eqform:status}
\end{table}

\section{Shared notation, conventions, and topology}
\label{sec:eqform:shared}

\subsection{Ensembles, amounts, volumes, and domains}

Unless a family states otherwise, temperature \(T>0\), imposed pressure
\(P^0>0\), and conserved inventory are fixed. Components are indexed by
\(i\in\mathcal C=\{1,\ldots,C\}\), candidate phases by
\(\alpha\in\mathcal P=\{1,\ldots,\Pi\}\), and independent reactions by
\(r\in\mathcal R=\{1,\ldots,R\}\). Phase amounts, totals, compositions, and
volumes are
\begin{equation}
 n_{i\alpha}\geq 0,\qquad
 n_\alpha=\sum_{i=1}^{C}n_{i\alpha},\qquad
 x_{i\alpha}=\frac{n_{i\alpha}}{n_\alpha},\qquad
 V_\alpha>0,\qquad v_\alpha=\frac{V_\alpha}{n_\alpha}.
\label{eq:eqform:phase-definitions}
\end{equation}
Composition is defined only when \(n_\alpha>0\). An \emph{active} phase has
\(n_\alpha>0\); an absent phase lies on \(n_\alpha=0\), where amount--volume
coordinates or a limiting tangent-plane construction must be used instead of
evaluating \(x_{i\alpha}\) by division. This is the common
appearance/disappearance topology.

For nonreactive problems the material balances are
\begin{equation}
 \sum_{\alpha\in\mathcal P} n_{i\alpha}=N_i,
 \qquad i\in\mathcal C.
\label{eq:eqform:material-balances}
\end{equation}
For reactive problems, let \(B\in\R^{E\times C}\) be a full-row-rank
element--species matrix and \(\boldsymbol b\) the conserved elemental
inventory. The independent conservation statement is
\begin{equation}
 B\sum_{\alpha\in\mathcal P}\boldsymbol n_\alpha=\boldsymbol b.
\label{eq:eqform:element-balances}
\end{equation}
Writing both Eq.~\eqref{eq:eqform:material-balances} for a fixed feed and
Eq.~\eqref{eq:eqform:element-balances} as independent constraints in a reactive
problem generally overconstrains the system. A reaction-extent formulation uses
\begin{equation}
 \boldsymbol n=\boldsymbol n^0+\nu\boldsymbol\xi,\qquad
 B\nu=0,\qquad \rank(\nu)=R.
\label{eq:eqform:reaction-extents}
\end{equation}
where the columns of \(\nu\in\R^{C\times R}\) are independent reaction vectors.
An extent representation spans every elemental-feasible composition direction
only when
\begin{equation}
 \operatorname{im}(\nu)=\ker(B),\qquad
 R=C-\rank(B).
\label{eq:eqform:reaction-completeness}
\end{equation}
If Eq.~\eqref{eq:eqform:reaction-completeness} is not satisfied, the extents
define a deliberately network-restricted feasible set
\(\boldsymbol n^0+\operatorname{im}(\nu)\), not an equivalent coordinate
system for elemental Gibbs minimization. Charge is included as a row of \(B\)
when it is an independent conserved quantity.

\subsection{Helmholtz evaluator and derivative contract}

For an active phase, the EOS supplies an extensive Helmholtz free energy
\begin{equation}
 A_\alpha=A(T,V_\alpha,\boldsymbol n_\alpha;\boldsymbol s_\alpha),
\label{eq:eqform:helmholtz-owner}
\end{equation}
where \(\boldsymbol s_\alpha\) is absent for a nonassociating evaluator and is
an implicit association state otherwise. Thermodynamic consistency requires
\begin{equation}
 P_\alpha=-\left(\frac{\partial A}{\partial V_\alpha}\right)_
 {T,\boldsymbol n_\alpha},\qquad
 \mu_{i\alpha}=\left(\frac{\partial A}{\partial n_{i\alpha}}\right)_
 {T,V_\alpha,n_{j\neq i,\alpha}}.
\label{eq:eqform:eos-derivatives}
\end{equation}
The isothermal--isobaric phase contribution is
\begin{equation}
 \mathcal G_\alpha(T,P^0,V_\alpha,\boldsymbol n_\alpha)
 =A(T,V_\alpha,\boldsymbol n_\alpha)+P^0V_\alpha.
\label{eq:eqform:phase-gibbs-transform}
\end{equation}
Its volume stationarity gives \(P_\alpha=P^0\), and its amount gradient gives
\(\boldsymbol\mu_\alpha\). These are identities in physical energy units; an
implementation may divide the complete objective and its derivatives by
\(RT\), but must not mix dimensional and dimensionless residuals.

At minimum, a direct-minimization formulation needs finite \(A\), \(P\), and
\(\boldsymbol\mu\) throughout its declared domain, together with gradients of
all constraints. A Newton or exact-Hessian NLP additionally needs consistent
second derivatives of the reduced objective. A formulation that constrains
chemical-potential equality needs derivatives of \(\boldsymbol\mu\), and an
exact Hessian of that residual problem can require third derivatives of \(A\).
Derivative availability is not a global-minimum certificate.

\subsection{Ions, electroneutrality, and the Galvani convention}

Let \(z_i\) be signed charge number and \(F\) Faraday's constant. Bulk phases
are locally electroneutral unless a formulation explicitly models interfacial
charge:
\begin{equation}
 \sum_{i=1}^{C}z_i n_{i\alpha}=0,\qquad \alpha\in\mathcal P.
\label{eq:eqform:phase-electroneutrality}
\end{equation}
With phase electric potential \(\Phi_\alpha\), the electrochemical potential is
\begin{equation}
 \widetilde\mu_{i\alpha}=\mu_{i\alpha}+z_iF\Phi_\alpha.
\label{eq:eqform:electrochemical-potential}
\end{equation}
\Status{verified thermodynamic convention}
Individual ionic chemical potentials are convention dependent. Therefore
individual ionic chemical-potential equality is not a certificate unless an
explicit common Galvani convention fixes every \(\Phi_\alpha\). The
electrolyte families below instead use charge-neutral modified-potential or
counterion-pair combinations in which the Galvani contribution cancels.

\subsection{Common KKT structure and phase appearance}

For a fixed candidate phase set in a neutral nonreactive problem, define
\begin{equation}
 G_{\mathcal P}^{*}=
 \min_{\{\boldsymbol n_\alpha,V_\alpha\}}
 \sum_{\alpha\in\mathcal P}\mathcal G_\alpha
 \quad\text{subject to Eq.~\eqref{eq:eqform:material-balances}, }
 \boldsymbol n_\alpha\geq0,\ V_\alpha>0.
\label{eq:eqform:generic-neutral-min}
\end{equation}
With material-balance multipliers \(\lambda_i\) and lower-bound multipliers
\(\gamma_{i\alpha}\geq0\), the amount KKT conditions are
\begin{equation}
 \mu_{i\alpha}-\lambda_i-\gamma_{i\alpha}=0,\qquad
 n_{i\alpha}\gamma_{i\alpha}=0,\qquad
 P_\alpha=P^0.
\label{eq:eqform:generic-neutral-kkt}
\end{equation}
Thus an interior component has equal chemical potential across active phases.
At zero amount, the inequality and complementarity condition replaces
equality. At phase disappearance, nonnegative tangent-plane distance for all
admissible trial phases supplies the topology certificate. Stationarity at
reported phases alone does not exclude an unreported lower-energy phase.

\paragraph{Acceptance template.}
A formulation-specific record must report: finite-domain and positivity
checks; independent balance and electroneutrality norms; objective or residual
convergence; stationarity/KKT or complementarity; an independent equilibrium
certificate; and the evidence, if any, for global phase-set completeness.
Tolerances are named inputs, not universal constants silently shared across
families.

\section{Public pressure-boundary and pure-saturation formulations}
\label{sec:eqform:public-boundary}

\subsection{Source, ensemble, and independent variables}

\Status{inferred thermodynamic specialization; verified implementation evidence}
The displayed phase-equilibrium conditions are derived from the verified
fixed-\(T,P\) Gibbs minimization/KKT construction in Pereira
et al.~\cite{Pereira2012} and agree with the direct constrained implementation
owner; Pereira does not specify these bubble/dew controllers. The retained
public routes are local
two-phase boundary problems, not executions of Pereira's HELD discovery
controller. For bubble pressure, \(T\) and parent-liquid composition
\(\boldsymbol x^{L,0}\) are fixed; unknowns are
\((P,\boldsymbol y,v_L,v_V)\), with \(C\) ambient composition entries and one
normalization constraint, hence \(C-1\) composition degrees of freedom. For
dew pressure, \(T\) and parent-vapour composition
\(\boldsymbol y^{V,0}\) are fixed; unknowns are
\((P,\boldsymbol x,v_L,v_V)\).

\subsection{Residual maps, domains, and topology}

The bubble-pressure residual map is
\begin{equation}
 \boldsymbol R_{\mathrm{bub}}=
 \begin{bmatrix}
 P_L(T,v_L,\boldsymbol x^{L,0})-P\\
 P_V(T,v_V,\boldsymbol y)-P\\
 \boldsymbol\mu_L(T,v_L,\boldsymbol x^{L,0})
 -\boldsymbol\mu_V(T,v_V,\boldsymbol y)\\
 \sum_i y_i-1
 \end{bmatrix}
 =\boldsymbol0.
\label{eq:eqform:bubble-residual}
\end{equation}
The dew-pressure residual interchanges fixed and trial compositions:
\begin{equation}
 \boldsymbol R_{\mathrm{dew}}=
 \begin{bmatrix}
 P_L(T,v_L,\boldsymbol x)-P\\
 P_V(T,v_V,\boldsymbol y^{V,0})-P\\
 \boldsymbol\mu_L(T,v_L,\boldsymbol x)
 -\boldsymbol\mu_V(T,v_V,\boldsymbol y^{V,0})\\
 \sum_i x_i-1
 \end{bmatrix}
 =\boldsymbol0.
\label{eq:eqform:dew-residual}
\end{equation}
Domains are \(P>0\), \(v_L>0\), \(v_V>v_L\), and compositions in the
simplex interior or on explicitly supported trace bounds. At an incipient
boundary the new phase amount tends to zero, so there is no finite two-phase
lever-rule material balance to impose. The fixed parent composition and unit
normalization of each phase define the boundary limit.

For one component, both composition vectors equal \(1\), leaving
\begin{equation}
 P(T,v_L)=P_{\mathrm{sat}},\qquad
 P(T,v_V)=P_{\mathrm{sat}},\qquad
 \mu(T,v_L)=\mu(T,v_V),\qquad v_V>v_L.
\label{eq:eqform:pure-saturation}
\end{equation}
This is the scoped public \texttt{single\_component\_vle} formulation.

\subsection{Objective/KKT interpretation and derivative contract}

\Evidence{verified implementation}
The current native owner represents both unit-normalized phases in amount and
volume variables, adds \(P\) as a route scalar, minimizes the dimensionless sum
of Eq.~\eqref{eq:eqform:phase-gibbs-transform}, and constrains fixed
composition, each phase total, both pressure residuals, all neutral
chemical-potential differences, and a positive phase-volume gap. That local
NLP implements Eqs.~\eqref{eq:eqform:bubble-residual}--%
\eqref{eq:eqform:pure-saturation}; it is not the global multiphase objective in
Eq.~\eqref{eq:eqform:generic-neutral-min}.

The evaluator must supply phase \(A,P,\boldsymbol\mu\); the residual Jacobian
needs first derivatives of \(P\) and \(\boldsymbol\mu\). The exposed exact
Hessian route needs corresponding higher derivatives. For narrow associating
bubble/dew proof scope, derivatives include the implicit association response.
The public pure route rejects association.

\subsection{Initialization, certificate, acceptance, and failure}

Density/volume roots must start in separated liquid and vapour regions and
remain inside EOS and volume bounds. Pressure and residual scaling must be
reported. The route-owned acceptance checks include pressure and neutral
chemical-potential consistency, fixed-parent composition, unit phase totals,
ln-fugacity consistency for non-pure boundaries, and the required
volume-gap/density-separation condition. These are local boundary checks, not a
stability certificate; no TPD stability certificate is currently public.
Multiple roots, root swapping, critical coalescence, trace-bound activity,
EOS-domain failure, and a vanishing phase-volume gap are explicit failures.
Ions would require
Eq.~\eqref{eq:eqform:electrochemical-potential} or a charge-neutral reduced
certificate and are outside this public family.

\paragraph{Relationship to other families.}
Bubble/dew and pure saturation solve one prescribed incipient topology. They do
not discover the number of phases, do not execute HELD Stages I--III, and do
not solve CE or CPE.

\section{Neutral HELD}
\label{sec:eqform:neutral-held}

\subsection{Source, scope, and canonical Stage III problem}

\Status{verified source formulation}
Pereira et al.~\cite{Pereira2012} define HELD for nonreactive multicomponent,
multiphase equilibrium at fixed \(T^0,P^0,\boldsymbol N\) using generic EOS
evaluations. Stage I searches for unstable trial states, Stage II alternates a
linear dual outer problem with nonconvex composition--volume lower problems,
and Stage III refines the candidate phase simplex.

For candidate set \(\mathcal P^*\), the canonical Stage III problem is direct
total-free-energy minimization:
\begin{equation}
 G_{\mathrm{HELD}}^{*}=
 \min_{\{\boldsymbol n_\alpha,V_\alpha\}_{\alpha\in\mathcal P^*}}
 \sum_{\alpha\in\mathcal P^*}
 \left[A(T^0,V_\alpha,\boldsymbol n_\alpha)+P^0V_\alpha\right],
\label{eq:eqform:held-stage-three}
\end{equation}
subject to
\begin{equation}
 \sum_{\alpha\in\mathcal P^*}n_{i\alpha}=N_i,\quad
 \boldsymbol n_\alpha^{L}\leq\boldsymbol n_\alpha\leq
 \boldsymbol n_\alpha^{U},\quad
 V_\alpha^{L}\leq V_\alpha\leq V_\alpha^{U}.
\label{eq:eqform:held-stage-three-constraints}
\end{equation}
The source uses a one-mole mass-number basis and tight candidate-neighbourhood
bounds. If the candidate set cannot satisfy balances, the problem is
infeasible and control returns to Stage II. The canonical Stage III problem is
direct total-free-energy minimization. Residual solving is not the canonical
Stage III formulation; it may be used only as a correction or diagnostic after
the direct minimization.

\subsection{Stages I and II, phase topology, and globality}

For normalized reference state \(\boldsymbol z\), a convenient
source-consistent tangent-plane representation is
\begin{equation}
 d(\boldsymbol x,v;\boldsymbol z)=
 g(T^0,P^0,\boldsymbol x,v)
 -g(T^0,P^0,\boldsymbol z,v_z)
 -\sum_{i=1}^{C-1}
 \left.\frac{\partial g}{\partial x_i}\right|_{\boldsymbol z}
 (x_i-z_i),
\label{eq:eqform:held-tpd}
\end{equation}
with volume included among lower-level variables and pressure stationarity
enforced at a minimizer. \Evidence{inference} Equation
\eqref{eq:eqform:held-tpd} fixes normalized notation for the source's
tangent-plane construction; the PDF's primal/dual equations remain the
authoritative source coordinates. A negative minimum identifies a missing
stable candidate. Stage II constructs upper and lower information from sampled
supporting planes. Finite multistart sampling does not prove the global minimum
of the nonconvex lower problem.

Phase appearance changes the candidate simplex. Duplicate candidates that
converge to the same composition and volume are merged. Trace species at source
lower bounds are refined in logarithmic coordinates. A candidate phase may
disappear at an active amount bound; complementarity and nonnegative TPD then
replace an undefined composition condition.

\subsection{KKT conditions and independent certificate}

For variables not active at candidate-neighbourhood bounds, the KKT equations
reduce to Eq.~\eqref{eq:eqform:generic-neutral-kkt}. Hence
\begin{equation}
 P_\alpha=P^0,\qquad
 \mu_{i\alpha}=\lambda_i
 \quad\text{for every interior component in every active phase}.
\label{eq:eqform:held-interior-kkt}
\end{equation}
\Evidence{inference} These relations follow from
Eqs.~\eqref{eq:eqform:held-stage-three}--%
\eqref{eq:eqform:held-stage-three-constraints}; they do not replace global dual
evidence.

The source's independent convergence certificate is
\begin{equation}
 0\leq UBD^{V}-G_{\mathrm{HELD}}^{*}\leq\epsilon_g,\qquad
 \max_{i,\alpha}
 \left|\frac{\mu_{i\alpha}-\mu_{i,\alpha+1}}{\mu_{i\alpha}}\right|
 \leq\epsilon_\mu,
\label{eq:eqform:held-certificate}
\end{equation}
together with feasibility, distinct-phase filtering, and trace refinement.
The chemical-potential check is independent of objective convergence but does
not alone prove global phase-set completeness. A stronger audit also requires
nonnegative globally minimized TPD over the admissible domain.

\subsection{Scaling, initialization, acceptance, and failure}

\Status{verified source numerics}
Pereira reports packing-fraction volume scaling, a one-mole mass basis, up to
\(10C\) Stage-I tunnelling starts, and Stage-II starts drawn from random,
shifted-previous, and near-pure states. Reported source values are
\(\epsilon_\lambda=0.5\), \(\epsilon_b=10^{-2}\),
\(\epsilon_x=\epsilon_\eta=10^{-3}\), and
\(\epsilon_\mu=\epsilon_g=10^{-6}\), with composition bounds
\(10^{-8}\) and \(1-10^{-8}\). These are historical source settings, not
invented repository defaults.

Acceptance requires a declared EOS domain, balance feasibility, converged
direct Stage III objective, KKT residuals, Eq.~\eqref{eq:eqform:held-certificate},
and an honest globality receipt for Stage-I/II lower searches. Failures include
missing candidate phases, infeasible Stage III balances, metastable local
minima, duplicate phases, active trace bounds, EOS failure, ill-scaled
mass/volume variables, and unproven lower-level globality.

\paragraph{Implementation and relationship.}
\Status{verified implementation evidence}
The repository contains neutral local NLP, deterministic candidate screening,
continuous TPD minimization, sampled-candidate audit, and local certification
components. It does not contain a completed Pereira Stage-I/II/III controller.
Neutral HELD is neither a direct boundary route nor CE/CPE.

\section{Association-aware neutral HELD}
\label{sec:eqform:association-held}

\subsection{Scope and implicit association state}

\Status{verified source scope; inferred derivative contract}
Pereira et al.~\cite{Pereira2012} use the same HELD controller with
association-capable SAFT evaluators and solve association equations at each EOS
evaluation. The controller and Stage III objective remain
Eqs.~\eqref{eq:eqform:held-stage-three}--%
\eqref{eq:eqform:held-stage-three-constraints}; this is an evaluator extension,
not a different phase-discovery algorithm.

Let association variables satisfy
\begin{equation}
 \boldsymbol R^{\mathrm{assoc}}
 (T,V,\boldsymbol n,\boldsymbol s)=\boldsymbol0,\qquad
 \boldsymbol s\in\mathcal S_{\mathrm{phys}},
\label{eq:eqform:association-state}
\end{equation}
and let \(\boldsymbol s^*(T,V,\boldsymbol n)\) be the unique physical solution
asserted by Pereira for this association system.
The reduced Helmholtz function used by HELD is
\begin{equation}
 \widehat A(T,V,\boldsymbol n)=
 A(T,V,\boldsymbol n;\boldsymbol s^*(T,V,\boldsymbol n)).
\label{eq:eqform:association-reduced-helmholtz}
\end{equation}
Every phase, trial TPD state, and derivative evaluation must solve
Eq.~\eqref{eq:eqform:association-state}; carrying a stale state between trial
points changes the mathematical function.

\subsection{Implicit derivatives}

If \(R_s=\partial\boldsymbol R^{\mathrm{assoc}}/\partial\boldsymbol s\)
is nonsingular at that solution, then for any outer variable \(y\),
\begin{equation}
 \frac{\dd\boldsymbol s^*}{\dd y}
 =-R_s^{-1}R_y,\qquad
 \frac{\dd\widehat A}{\dd y}
 =A_y+A_s\frac{\dd\boldsymbol s^*}{\dd y}.
\label{eq:eqform:association-implicit-derivative}
\end{equation}
\Evidence{inference} Equation
\eqref{eq:eqform:association-implicit-derivative} is the required
implicit-function derivative contract. Pereira does not print this equation,
the association residual, its tolerance, or the exact reduced-Hessian formula.
Those details are \Evidence{unknown} at primary-source level. A local NLP with
exact Hessians must differentiate
Eq.~\eqref{eq:eqform:association-implicit-derivative} again or use an
equivalent exact automatic/implicit construction.

\subsection{KKT, certificate, and numerical behavior}

The outer KKT and phase-topology conditions are
Eqs.~\eqref{eq:eqform:held-interior-kkt} and
\eqref{eq:eqform:held-certificate}, evaluated with \(\widehat A\) and its total
derivatives. \Evidence{inference} A designed repository acceptance certificate
adds
\begin{equation}
 \max_{\alpha\in\mathcal P^*}
 \norm{\boldsymbol R^{\mathrm{assoc}}_\alpha}_{\infty}
 \leq\epsilon_{\mathrm{assoc}},
\label{eq:eqform:association-certificate}
\end{equation}
 plus a derivative-consistency receipt; Pereira does not define this
\(\epsilon_{\mathrm{assoc}}\) test. Pereira states that its association system
has a unique solution and solves it independently at each EOS evaluation.
Broader multiple-root behavior for other association evaluators is
\Evidence{unknown}, not a Pereira result. The verified root can still become
ill-conditioned, \(R_s\) may approach singularity, inner tolerances can pollute
outer stationarity, and EOS calls can fail in Stage-I/II trial states.
Initialization must declare the inner-solve initialization, warm-start, and
continuation policy without silently changing the selected thermodynamic
branch.

\paragraph{Implementation and non-equivalence.}
\Status{verified implementation evidence}
The repository has association mass-action residual/Jacobian components,
implicit sensitivities, and exact derivative corrections. A generalized
association-aware HELD controller remains absent/deferred. Narrow public
associating bubble/dew evidence does not admit this family.

\section{Perdomo modified-mole HELD2}
\label{sec:eqform:perdomo-held2}

\subsection{Source, scope, and eliminated-ion coordinates}

\Status{verified source formulation}
Perdomo et al.~\cite{Perdomo2025} treat nonreactive strong electrolytes in
fluid phases at fixed \(T,P^0\), with at least one molecular species and
complete dissociation. Solids, weak-electrolyte reactions, and CPE are outside
source scope.

Choose charged species \(E\), with \(z_E\neq0\). Local electroneutrality
eliminates its amount:
\begin{equation}
 n_E=-\frac{1}{z_E}\sum_{i\neq E}z_i n_i.
\label{eq:eqform:perdomo-elimination}
\end{equation}
For \(i\in\mathcal C^{(E)}=\mathcal C\setminus\{E\}\), define
\begin{equation}
 \overline n_i^{(E)}
 =\left(1-\frac{z_i}{z_E}\right)n_i,\qquad
 \overline n=\sum_{i\in\mathcal C^{(E)}}\overline n_i^{(E)},\qquad
 \overline x_i^{(E)}=\frac{\overline n_i^{(E)}}{\overline n}.
\label{eq:eqform:perdomo-modified-moles}
\end{equation}
The modified fractions sum to one. There are \(C-1\) retained modified species
and \(C-2\) independent modified compositions; volume is the additional
lower-level variable. Choose one dependent modified species \(q\in
\mathcal C^{(E)}\), define
\(\mathcal C^{(EC)}=\mathcal C^{(E)}\setminus\{q\}\), and reconstruct
\(\overline x_q^{(E)}=1-\sum_{i\in\mathcal C^{(EC)}}
\overline x_i^{(E)}\). Thus
\(\overline{\boldsymbol x}^{(EC)}\in\R^{C-2}\) is the independent coordinate
vector used below.

Define charge-constrained and modified potentials by
\begin{equation}
 \mu_i^{(E)}
 =\widetilde\mu_i-\frac{z_i}{z_E}\widetilde\mu_E,\qquad
 \overline\mu_i^{(E)}
 =\frac{\mu_i^{(E)}}{1-z_i/z_E}.
\label{eq:eqform:perdomo-modified-potential}
\end{equation}
Galvani terms cancel in \(\mu_i^{(E)}\). The source's eliminated-species
representation can assign a phase-potential relation, but it must not be
promoted to a unique measurable Galvani potential.

\subsection{Canonical Stage III objective and constraints}

For \(m_p\) candidate phases, fractions \(\phi^\alpha\in[0,1]\), molar
volumes \(v_\alpha\), and modified compositions
\(\overline{\boldsymbol x}^{(EC),\alpha}\), Perdomo's direct Stage III is
\begin{equation}
 \overline G_{\mathrm{HELD2}}^{*}
 =\min_{\{\phi^\alpha,v_\alpha,\overline{\boldsymbol x}^{(E),\alpha}\}}
 \sum_{\alpha=1}^{m_p}\phi^\alpha
 \left[
 \overline A^{\mathrm{el}}
 (T,v_\alpha,\overline{\boldsymbol x}^{(EC),\alpha})
 +P^0v_\alpha
 \right],
\label{eq:eqform:perdomo-stage-three}
\end{equation}
subject to
\begin{equation}
 \sum_{\alpha=1}^{m_p}\phi^\alpha
 \overline x_i^{(E),\alpha}
 =\overline x_{0,i}^{(E)},\quad i\in\mathcal C^{(EC)},\qquad
 \sum_{\alpha=1}^{m_p}\phi^\alpha=1,\quad
 \phi^\alpha\geq0,\quad
 \overline{\boldsymbol x}^{(EC),\alpha}\in\overline{\mathcal X}^{(EC)},\quad
 v_\alpha\in[v^L,v^U].
\label{eq:eqform:perdomo-stage-three-constraints}
\end{equation}
Candidate-neighbourhood composition bounds are also imposed in the source.
The \(C-2\) displayed modified balances and the phase-fraction normalization
are independent; the omitted modified balance follows from modified-fraction
normalization. Together with Eq.~\eqref{eq:eqform:perdomo-elimination}, they
recover ordinary balances and local charge neutrality.

\subsection{KKT, certificate, scaling, and failure}

For an active phase and interior modified component, KKT conditions give
\begin{equation}
 P_\alpha=P^0,\qquad
 \overline\mu_i^{(E),\alpha}
 =\overline\lambda_i,
\label{eq:eqform:perdomo-kkt}
\end{equation}
so the independent interphase condition is equality of modified potentials,
not convention-dependent individual ionic potentials. The source certificate
is
\begin{equation}
 0\leq UBD^{V}-\overline G_{\mathrm{HELD2}}^{*}\leq\epsilon_g,\qquad
 \max_{i,\alpha}
 \left|
 \frac{\overline\mu_i^{(E),\alpha}
 -\overline\mu_i^{(E),\alpha+1}}
 {\overline\mu_i^{(E),\alpha}}
 \right|\leq\epsilon_\mu,
\label{eq:eqform:perdomo-certificate}
\end{equation}
plus modified-balance, normalization, volume, duplicate-phase, and
electroneutrality checks.

\Status{verified source numerics}
The eliminated ion is chosen with largest absolute charge. The source uses
fixed-seed \(10C\) stability starts, packing-fraction bounds
\(10^{-6}\leq\eta\leq0.74\), modified lower fractions initially proportional
to \(10^{-10}(1-z_i/z_E)\), logarithmic trace refinement as low as
\(10^{-300}\), and default reported
\(\epsilon_g=\epsilon_\mu=10^{-6}\). Nonconvex multistart lower solves do not
provide a deterministic global guarantee. Exact derivative provenance and
Hessian use are \Evidence{unknown}.

Failures include a nonpositive eliminated amount, singular modified coordinate
factor, infeasible modified balances, missing phase candidates, trace-bound
distortion, metastable lower solves, and inconsistent recovery to ordinary ion
amounts.

\paragraph{Implementation and non-equivalence.}
\Status{verified implementation evidence}
No source-faithful modified-mole HELD2 controller is implemented. Historical
repository names containing \texttt{held2} operate on counterion-pair
coordinates and do not establish Perdomo parity. Perdomo and Ascani
formulations are not equated.

\section{Ascani counterion-pair equilibrium}
\label{sec:eqform:ascani-pairs}

\subsection{Source, coordinates, and residual map}

\Status{verified source formulation}
Ascani, Sadowski, and Held~\cite{Ascani2022} treat fixed-\(T,P\)
multiphase electrolyte equilibrium with locally electroneutral bulk phases.
Their published controller is a counterion-pair residual and successive
phase-addition algorithm, not HELD2.

Let \(\mathcal C_{\mathrm{ch}}\) contain \(C_{\mathrm{ch}}\) charged species.
Construct \(C_{\mathrm{ch}}-1\) independent cation--anion pairs and
\(E\in\R^{(C_{\mathrm{ch}}-1)\times C_{\mathrm{ch}}}\), with
\begin{equation}
 \rank(E)=C_{\mathrm{ch}}-1,\qquad
 E\boldsymbol z_{\mathrm{ch}}=\boldsymbol0.
\label{eq:eqform:ascani-pair-rank}
\end{equation}
A row has inverse-absolute-valence weights on its selected cation and anion.
The independent pair potential is
\begin{equation}
 \widehat{\boldsymbol\mu}_{\pm,\alpha}
 =E\widetilde{\boldsymbol\mu}_{\mathrm{ch},\alpha}
 =E\boldsymbol\mu_{\mathrm{ch},\alpha}.
\label{eq:eqform:ascani-pair-potential}
\end{equation}
The second equality follows from Eq.~\eqref{eq:eqform:ascani-pair-rank}; the
Galvani term cancels. Row normalization changes numerical scale but not
cross-phase equality.

Let \(\mathcal S_{\mathrm{ch}}\) select \(C_{\mathrm{ch}}-1\) independent
charged-species balances. Taking phase \(1\) as reference, the source's
\emph{reduced} equilibrium residual is
\begin{equation}
 \boldsymbol R_{\mathrm{Ascani}}=
 \begin{bmatrix}
 \{\mu_{i\alpha}-\mu_{i1}\}_{z_i=0,\ \alpha>1}\\
 \{\widehat{\boldsymbol\mu}_{\pm,\alpha}
 -\widehat{\boldsymbol\mu}_{\pm,1}\}_{\alpha>1}\\
 \{\sum_\alpha n_{i\alpha}-N_i\}_{z_i=0}\\
 \{\sum_\alpha n_{k\alpha}-N_k\}_{k\in\mathcal S_{\mathrm{ch}}}
 \end{bmatrix}
 =\boldsymbol0,
\label{eq:eqform:ascani-residual}
\end{equation}
written in reduced pair-transfer variables that move selected ions in
inverse-valence proportions and preserve phase electroneutrality. The omitted
charged balance follows from globally neutral feed plus per-phase neutrality.
No explicit charge row appears because neutrality is an identity of the
coordinates. This is source Eqs.~(22), (23), (31), and (32), not a mixture of
that reduced system with the full Eqs.~(22)--(25).

\subsection{Stability, topology, KKT interpretation, and certificate}

The source defines transformed tangent-plane distance over the same
electroneutral pair coordinates:
\begin{equation}
 \TPD_{\mathrm{pair}}(\boldsymbol\xi)
 =g(\boldsymbol n(\boldsymbol\xi))
 -g(\boldsymbol n^0)
 -\boldsymbol\mu_{\mathrm{neutral}}^{0\mathsf T}
 \Delta\boldsymbol n_{\mathrm{neutral}}
 -\widehat{\boldsymbol\mu}_{\pm}^{0\mathsf T}
 \Delta\boldsymbol n_{\mathrm{pair}}.
\label{eq:eqform:ascani-tpd}
\end{equation}
\Evidence{inference} Equation \eqref{eq:eqform:ascani-tpd} is a compact
coordinate statement of source Eqs.~(33)--(35); the source transform defines
the exact \(\boldsymbol n(\boldsymbol\xi)\).

A negative TPD trial generates a split, after which a modified-Powell solve of
Eq.~\eqref{eq:eqform:ascani-residual} is attempted. Stability testing and phase
addition repeat. KKT interpretation is equality of neutral potentials and of
the \(C_{\mathrm{ch}}-1\) independent pair potentials, with balances and local
charge constraints. The independent certificate is small residual norm, full
pair-matrix rank, and balance and charge feasibility. \Evidence{verified} The
published controller retests the light phase. \Evidence{inference} Requiring
nonnegative admissible pair TPD for every retained phase and objective
comparison across phase sets is a stronger designed certificate needed before
a globality claim; its completeness is not established by Ascani.

The source reports random search and shrinking perturbations but does not give
a deterministic seed, complete tolerance set, Jacobian construction, or formal
global guarantee; these are \Evidence{unknown}. Failures include a
rank-deficient pair basis, infeasible reduced coordinates, negative amounts,
missed negative-TPD regions, local residual roots, duplicate phases, and
ill-scaled multivalent pairs.

\paragraph{Implementation and non-equivalence.}
\Status{verified implementation evidence}
Counterion-pair matrices, reduced coordinates, mean-ionic bookkeeping, TPD and
postsolve components exist internally. The full source controller remains
absent/deferred. This is separate from Perdomo's eliminated-ion modified-mole
HELD2 and from reactive CE/CPE.

\section{Standalone chemical equilibrium}
\label{sec:eqform:ce}

\subsection{Ensemble, objective, constraints, and domains}

\Status{verified source basis}
Specializing the thermodynamic construction in Ascani~\cite{Ascani2023} to one
homogeneous phase gives standalone CE at fixed \(T,P^0,\boldsymbol b\):
\begin{equation}
 G_{\mathrm{CE}}^{*}
 =\min_{\boldsymbol n\geq0}
 G(T,P^0,\boldsymbol n)
 \quad\text{subject to}\quad
 B\boldsymbol n=\boldsymbol b.
\label{eq:eqform:ce-objective}
\end{equation}
If charged species are present, bulk electroneutrality is an additional
independent constraint unless charge conservation is already an independent
row of \(B\). Standard states and activity conventions must be declared; no
single-ion standard chemical potential is observable without a convention.

When Eq.~\eqref{eq:eqform:reaction-completeness} holds, extent coordinates are
equivalent and Eq.~\eqref{eq:eqform:reaction-extents} gives
\begin{equation}
 \boldsymbol R_{\mathrm{CE}}(\boldsymbol\xi)
 =\nu^\mathsf T\boldsymbol\mu
 (T,P^0,\boldsymbol n^0+\nu\boldsymbol\xi)
 =\boldsymbol0,
\label{eq:eqform:ce-affinity}
\end{equation}
on \(\boldsymbol n^0+\nu\boldsymbol\xi\geq0\). Under a declared activity
standard state, each row is equivalently
\begin{equation}
 \sum_{i=1}^{C}\nu_{ir}\mu_i
 =\Delta_r G_r^0+RT\sum_{i=1}^{C}\nu_{ir}\ln a_i=0.
\label{eq:eqform:ce-mass-action}
\end{equation}
For a deliberately specified reaction network that does not span \(\ker B\),
Eqs.~\eqref{eq:eqform:ce-affinity}--\eqref{eq:eqform:ce-mass-action} certify
only the restricted extent problem
\(\min_{\boldsymbol\xi}G(T,P^0,\boldsymbol n^0+\nu\boldsymbol\xi)\), not the
elemental problem in Eq.~\eqref{eq:eqform:ce-objective}.

\subsection{KKT, boundary species, and certificate}

With multipliers \(\boldsymbol\lambda\) for \(B\boldsymbol n=\boldsymbol b\)
and \(\boldsymbol\gamma\geq0\) for nonnegativity,
\begin{equation}
 \boldsymbol\mu-B^\mathsf T\boldsymbol\lambda-\boldsymbol\gamma
 =\boldsymbol0,\qquad
 n_i\gamma_i=0.
\label{eq:eqform:ce-kkt}
\end{equation}
For a complete interior extent basis, Eq.~\eqref{eq:eqform:ce-kkt} implies
Eq.~\eqref{eq:eqform:ce-affinity}. At zero species amount, complementarity
replaces invalid logarithmic equality.

An independent CE certificate reports elemental and charge balances,
positivity/complementarity, affinities, standard-state identity, and objective
comparison. If the reduced Gibbs objective is proven convex over the feasible
set, KKT is sufficient globally; otherwise an affinity root is only necessary.
Initialization must use a strictly feasible composition or a declared boundary
strategy. An incomplete undeclared reaction span, rank-deficient \(B\) or
\(\nu\), dependent reactions, inconsistent
standard states, log-domain failure, boundary species, multiple stationary
points, and nonconvex activities are explicit failures.

\paragraph{Topology, implementation, and relationship.}
Standalone CE has one homogeneous phase and no phase-appearance topology or
interphase certificate. \Status{verified implementation evidence}
Public schemas and helpers exist, while the executable CE solver is internal
validation only and the retained nonideal proof is not accepted. CE is not
HELD2, counterion-pair phase equilibrium, or CPE.

\section{Coupled phase-and-chemical equilibrium}
\label{sec:eqform:cpe}

\subsection{Thermodynamic problem and independent constraints}

\Status{verified source Gibbs problem; inferred Helmholtz specialization;
separately designed repository family}
For fixed \(T,P^0,\boldsymbol b\), simultaneous CPE minimizes over phase count,
phase amounts, and phase volumes:
\begin{equation}
 G_{\mathrm{CPE}}^{*}
 =\min_{\mathcal P,\{\boldsymbol n_\alpha,V_\alpha\}}
 \sum_{\alpha\in\mathcal P}
 \left[A(T,V_\alpha,\boldsymbol n_\alpha)+P^0V_\alpha\right],
\label{eq:eqform:cpe-objective}
\end{equation}
subject to
\begin{equation}
 B\sum_{\alpha\in\mathcal P}\boldsymbol n_\alpha=\boldsymbol b,\qquad
 \boldsymbol n_\alpha\geq0,\quad V_\alpha>0,
\label{eq:eqform:cpe-constraints}
\end{equation}
and, when ions occur, independent per-phase electroneutrality constraints.
Ascani~\cite{Ascani2023} verifies the total-Gibbs problem; rewriting each
phase contribution as \(A+P^0V\) in
Eq.~\eqref{eq:eqform:cpe-objective} is an \Evidence{inference} from the shared
Legendre transform. A stoichiometric version may introduce a complete
reaction-extent basis and a reference-phase balance. A network-restricted
version must say so explicitly, and neither version may duplicate elemental
constraints.

\subsection{KKT/residual conditions and phase topology}

For neutral interior species, derived KKT conditions imply
\begin{equation}
 P_\alpha=P^0,\qquad
 \mu_{i\alpha}=\mu_{i\beta},\qquad
 \nu^\mathsf T\boldsymbol\mu_\alpha=\boldsymbol0,
\label{eq:eqform:cpe-kkt}
\end{equation}
for active phases \(\alpha,\beta\), with inequality/complementarity at zero
amounts. The affinity rows cover every elemental-feasible reaction direction
only under Eq.~\eqref{eq:eqform:reaction-completeness}; otherwise they certify
only the declared reaction network. If ions are present, interphase conditions
must use explicitly
Galvani-referenced electrochemical potentials or an independent charge-neutral
modified/pair basis; generic individual ionic equality is forbidden.

Ascani's demonstrated final solver uses coupled interphase-fugacity and
reaction-affinity residuals after homogeneous CE, nonreactive TPD, and
phasewise CE initialization. Those preliminary solves initialize CPE; they are
not CPE themselves. The paper explicitly notes that interphase and reaction
residual roots are necessary but not sufficient and may include trivial or
metastable solutions.

Phase appearance/disappearance combines topology of
Eq.~\eqref{eq:eqform:generic-neutral-kkt} with reactive feasible directions.
An absent phase must have nonnegative \emph{reactive} tangent-plane distance
over compositions reachable under elemental conservation; nonreactive TPD
alone is not a complete certificate.

\subsection{Certificate, numerics, and failure}

A complete CPE certificate must independently establish elemental and charge
feasibility, pressure stationarity, interphase neutral or charge-neutral
potential equality, reaction affinities, complementarity, nonnegative reactive
phase-appearance tests, and objective comparison across competing phase counts
and reaction branches. \Evidence{unknown} The Ascani source supplies no final
reactive global-minimum certificate, reactive-TPD completeness proof, or
electrolyte/Galvani extension. Those requirements must be designed and tested
before capability admission.

Scaling must distinguish energy, reaction-affinity, material, volume, and
charge rows. Initialization may use the source's staged CE/TPD/phasewise-CE
sequence, but final acceptance must depend on simultaneous solve and its
independent certificate. Failures include inconsistent elemental inventory,
dependent reaction bases, missed phases, nonreactive-stability false positives,
trivial or metastable roots, phase extinction, boundary species, and
incompatible standard states.

\paragraph{Implementation and non-equivalence.}
\Status{verified implementation evidence}
No runtime CPE formulation is implemented. CPE is not standalone CE followed
by a nonreactive flash, and it is not silently identified with CE, HELD2, or
counterion-pair equilibrium.

\section{Source-to-formulation reconciliation}
\label{sec:eqform:reconciliation}

\small
\begin{longtable}{>{\raggedright\arraybackslash}p{0.13\textwidth}
                  >{\raggedright\arraybackslash}p{0.24\textwidth}
                  >{\raggedright\arraybackslash}p{0.25\textwidth}
                  >{\raggedright\arraybackslash}p{0.27\textwidth}}
\caption{Source reconciliation. Owner means equation owner in this document,
not public runtime ownership.}
\label{tab:eqform:reconciliation}\\
\toprule
Owner & Source equation or statement & Evidence and source anchor &
Repository formulation owner / unresolved point \\
\midrule
\endfirsthead
\toprule
Owner & Source equation or statement & Evidence and source anchor &
Repository formulation owner / unresolved point \\
\midrule
\endhead
Eqs.~\eqref{eq:eqform:phase-definitions}--%
\eqref{eq:eqform:phase-gibbs-transform} &
Extensive Helmholtz/Gibbs variables and derivatives &
\Evidence{verified}: Pereira 2012, Sections 2--3; current phase evaluator &
Shared notation only; EOS equations remain in their separate canonical owner. \\
\addlinespace
Eqs.~\eqref{eq:eqform:phase-electroneutrality}--%
\eqref{eq:eqform:electrochemical-potential} &
Local charge and electrochemical potential &
\Evidence{verified}: Ascani 2022, Eqs.~(5), (9)--(15) &
Shared ionic convention; absolute Galvani potential remains convention
dependent. \\
\addlinespace
Eqs.~\eqref{eq:eqform:bubble-residual}--%
\eqref{eq:eqform:pure-saturation} &
Incipient two-phase pressure and chemical-potential conditions &
\Evidence{inference}: specialization of verified fixed-\(T,P\) KKT in Pereira
2012; residual structure and pure NIST scope are verified implementation
evidence &
Public direct boundary formulation. A paper-specific global boundary
controller is \Evidence{unknown} and not claimed. \\
\addlinespace
Eqs.~\eqref{eq:eqform:held-stage-three}--%
\eqref{eq:eqform:held-stage-three-constraints} &
HELD Stage III direct total-free-energy problem &
\Evidence{verified}: Pereira 2012, Section 3.3.1, problem
\(\mathrm{GM}_{x,V}\) &
Neutral HELD owner; controller absent/deferred. \\
\addlinespace
Eq.~\eqref{eq:eqform:held-tpd} &
Normalized TPD notation &
\Evidence{inference}: compact rewrite of Pereira 2012, Eqs.~(1)--(8) &
Exact PDF primal/dual coordinates remain authoritative. \\
\addlinespace
Eqs.~\eqref{eq:eqform:held-interior-kkt}--%
\eqref{eq:eqform:held-certificate} &
Interior stationarity, energy gap, chemical-potential convergence &
\Evidence{verified}: Pereira 2012, Eqs.~(12)--(13);
KKT reduction is \Evidence{inference} &
No deterministic global guarantee from finite multistart. \\
\addlinespace
Eqs.~\eqref{eq:eqform:association-state}--%
\eqref{eq:eqform:association-reduced-helmholtz} &
Association-aware evaluator and unique inner solution &
Inner association solve is \Evidence{verified}: Pereira 2012, Section 4;
the source states uniqueness &
Association residual details and primary-source tolerance are
\Evidence{unknown}. \\
\addlinespace
Eqs.~\eqref{eq:eqform:association-implicit-derivative}--%
\eqref{eq:eqform:association-certificate} &
Implicit derivative and repository acceptance contract &
\Evidence{inference}: implicit-function derivation and designed acceptance
test; not printed by Pereira &
Residual form, primary-source tolerance, and exact source Hessian are
\Evidence{unknown}. \\
\addlinespace
Eqs.~\eqref{eq:eqform:perdomo-elimination}--%
\eqref{eq:eqform:perdomo-modified-potential} &
Eliminated ion, modified moles/fractions/potentials &
\Evidence{verified}: Perdomo 2025, Eqs.~(9), (15), (23), and (28)--(31) &
Perdomo HELD2 owner; implementation absent/deferred. \\
\addlinespace
Eqs.~\eqref{eq:eqform:perdomo-stage-three}--%
\eqref{eq:eqform:perdomo-certificate} &
Modified-mole Stage III and certificate &
\Evidence{verified}: Perdomo 2025, Eqs.~(67)--(69) &
Derivative/Hessian provenance is \Evidence{unknown}. \\
\addlinespace
Eqs.~\eqref{eq:eqform:ascani-pair-rank}--%
\eqref{eq:eqform:ascani-residual} &
Independent pair basis and equilibrium residual &
\Evidence{verified}: Ascani 2022, Eqs.~(22), (23), and (26)--(32); the displayed
residual is the reduced system &
Pair components internal; complete controller absent/deferred. \\
\addlinespace
Eq.~\eqref{eq:eqform:ascani-tpd} &
Electroneutral transformed TPD &
\Evidence{inference}: compact rewrite of Ascani 2022, Eqs.~(33)--(35) &
Search counts, seed, tolerances, and global guarantee are
\Evidence{unknown}. \\
\addlinespace
Eqs.~\eqref{eq:eqform:ce-objective}--%
\eqref{eq:eqform:ce-mass-action} &
Homogeneous Gibbs minimization and reaction affinities &
\Evidence{verified}: Ascani 2023, Eqs.~(1)--(10), specialized to one phase;
reaction-space equivalence requires
Eq.~\eqref{eq:eqform:reaction-completeness} &
Rawlings--Ekerdt convexity details are retained only in a secondary local note;
primary verification is \Evidence{unknown}. \\
\addlinespace
Eq.~\eqref{eq:eqform:ce-kkt} &
Species nonnegativity complementarity &
\Evidence{inference}: KKT expansion of the verified Gibbs problem &
Boundary-species acceptance remains a repository design requirement. \\
\addlinespace
Eqs.~\eqref{eq:eqform:cpe-objective}--%
\eqref{eq:eqform:cpe-constraints} &
Simultaneous total-Gibbs CPE &
\Evidence{verified}: Ascani 2023, Eqs.~(1)--(10), for total Gibbs;
the displayed \(A+P^0V\) transform is \Evidence{inference} &
Elemental-inventory direct formulation is the canonical repository design. \\
\addlinespace
Eq.~\eqref{eq:eqform:cpe-kkt} &
Interphase and reaction necessary conditions plus complementarity &
\Evidence{verified}: Ascani 2023 for interphase/reaction necessary conditions;
explicit Helmholtz-coordinate KKT and complementarity are
\Evidence{inference} &
Stoichiometric-matrix orientation in source is dimensionally ambiguous; final
reactive global certificate is \Evidence{unknown}. \\
\bottomrule
\end{longtable}
\normalsize

\section{Thermodynamic and numerical consistency review}
\label{sec:eqform:consistency}

\small
\begin{longtable}{>{\raggedright\arraybackslash}p{0.15\textwidth}
                  >{\raggedright\arraybackslash}p{0.19\textwidth}
                  >{\raggedright\arraybackslash}p{0.18\textwidth}
                  >{\raggedright\arraybackslash}p{0.18\textwidth}
                  >{\raggedright\arraybackslash}p{0.17\textwidth}}
\caption{Internal consistency audit.}
\label{tab:eqform:consistency}\\
\toprule
Family & Degrees of freedom / rank & Units and basis &
Certificate independence & Phase-degeneracy behavior \\
\midrule
\endfirsthead
\toprule
Family & Degrees of freedom / rank & Units and basis &
Certificate independence & Phase-degeneracy behavior \\
\midrule
\endhead
Public boundaries &
Bubble/dew use \(C+3\) ambient unknowns: \(P\), two volumes, and all \(C\)
trial entries. Two pressure rows, \(C\) potential rows, and normalization give
\(C+3\) equations; normalization leaves \(C-1\) composition degrees. The pure
case leaves three unknowns and three independent equations. &
Pressure in Pa, molar volume in \(\mathrm{m^3\,mol^{-1}}\), potentials in
\(\mathrm{J\,mol^{-1}}\), or one declared complete \(RT\) scaling. &
Pressure, potential, normalization, and phase-gap checks are distinct from NLP
termination. &
At critical coalescence \(v_V-v_L\to0\); route must reject or change topology. \\
\addlinespace
Neutral HELD &
For fixed \(\Pi\), \(\Pi(C+1)\) amount/volume variables minus \(C\) independent
balances before active bounds. &
Pereira Stage III uses extensive free energy and a one-mole mass-number basis;
conversion must be explicit. &
Dual/free-energy gap and cross-phase potential checks are independent; global
TPD is stronger than either local check. &
Zero phase amount uses complementarity; duplicate phases are merged; infeasible
candidate simplex returns to Stage II. \\
\addlinespace
Association HELD &
Neutral HELD degrees plus association states eliminated by a square,
nonsingular \(R_s\). &
Association residual and energy derivatives must use one declared site and
amount basis. &
Association residual and derivative receipt are additional to HELD
certificates. &
Inner-root loss or singular \(R_s\) invalidates the phase evaluator before
outer topology decisions. \\
\addlinespace
Perdomo HELD2 &
\(C-2\) independent \(\mathcal C^{(EC)}\) coordinates and balances plus phase
normalization avoid the dependent modified-fraction row; volume adds one lower
variable. &
Modified moles preserve total mole basis under electroneutral elimination;
potentials retain energy-per-mole units. &
Modified-potential equality and dual/free-energy gap are independent; ordinary
individual-ion equality is not used. &
\(\phi^\alpha=0\) removes a phase; trace modified fractions require reversible
recovery to physical ion amounts. \\
\addlinespace
Ascani pairs &
Pair matrix rank is \(C_{\mathrm{ch}}-1\); charge-preserving variables remove
redundant neutrality directions. The reduced system retains
\(C_{\mathrm{ch}}-1\) ionic balances and no explicit charge rows. &
Pair row scaling must be declared; inverse-valence weighting preserves
energy-per-mole combinations up to row scale. &
Mean-ionic/pair equality, balance/charge, and TPD are separate checks. &
Successive phase addition handles appearance; missed TPD regions or duplicate
roots leave topology unproven. \\
\addlinespace
Standalone CE &
\(\rank(B)\) independent conservation rows and \(\rank(\nu)=R\) independent
reaction directions. Equivalence needs
\(\operatorname{im}(\nu)=\ker(B)\) and \(R=C-\rank(B)\); otherwise the network
is explicitly restricted. &
Activities are dimensionless and require declared standard states; affinities
are energy per reaction extent. &
Balances, complementarity, affinities, and objective/global-convexity evidence
are distinct. &
No phase topology; zero species amounts activate inequalities and invalidate
naive log residuals. \\
\addlinespace
CPE &
Use either full-row-rank elemental constraints or an independent extent/
reference-balance system, not both redundantly. Extents must be complete or
explicitly network restricted. &
All phases share one energy, amount, pressure, and standard-state convention. &
Coupled residual convergence is necessary; reactive TPD and global objective
comparison are additional and currently unknown. &
Appearance requires reactive feasible directions; disappearance combines
phase and species complementarity. \\
\bottomrule
\end{longtable}
\normalsize

\paragraph{Review conclusion.}
\Status{verified internal review}
The seven families use consistent fixed-\(T,P\) thermodynamic bases while
remaining non-equated. Balance ranks, amount/volume bases, ionic conventions,
derivative requirements, certificate independence, and phase-degeneracy
behavior are explicit. Unresolved items in
Table~\ref{tab:eqform:reconciliation} remain unknown rather than being filled
with parameters or defaults.

\begin{thebibliography}{9}

\bibitem{Pereira2012}
F.~E. Pereira, G.~Jackson, A.~Galindo, and C.~S. Adjiman,
\emph{The HELD algorithm for multicomponent, multiphase equilibrium
calculations with generic equations of state},
Computers \& Chemical Engineering 36 (2012) 99--118.
\href{https://doi.org/10.1016/j.compchemeng.2011.07.009}
{doi:10.1016/j.compchemeng.2011.07.009}.

\bibitem{Perdomo2025}
F.~A. Perdomo, G.~Jackson, A.~Mitsos, A.~Galindo, and C.~S. Adjiman,
\emph{Phase stability criteria and fluid-phase equilibria in
strong-electrolyte systems},
Computers \& Chemical Engineering 194 (2025) 108977.
\href{https://doi.org/10.1016/j.compchemeng.2024.108977}
{doi:10.1016/j.compchemeng.2024.108977}.

\bibitem{Ascani2022}
M.~Ascani, G.~Sadowski, and C.~Held,
\emph{Calculation of Multiphase Equilibria Containing Mixed Solvents and
Mixed Electrolytes: General Formulation and Case Studies},
Journal of Chemical \& Engineering Data 67 (2022) 1972--1984.
\href{https://doi.org/10.1021/acs.jced.1c00866}
{doi:10.1021/acs.jced.1c00866}.

\bibitem{Ascani2023}
M.~Ascani, G.~Sadowski, and C.~Held,
\emph{Simultaneous Predictions of Chemical and Phase Equilibria in Systems
with an Esterification Reaction Using PC-SAFT},
Molecules 28 (2023) 1768.
\href{https://doi.org/10.3390/molecules28041768}
{doi:10.3390/molecules28041768}.

\end{thebibliography}

\end{document}
